% PDF 266--275；原书第 238--247 页。
\subsection{缓解策略}
\label{sec:ch06-debiasing}
下面按 \refcite{ch06-ref-029}{Gallegos 等人（2023）}的论述，考察偏见缓解策略，并依其在大语言模型工作流程中所处的阶段，区分预处理、内部处理和后处理，如图~\ref{fig:ch06-debiasing} 所示。

\begin{figure}[!htbp]
\centering
% 重绘自原书第 238 页的图 6.3。
\begin{tikzpicture}[
  font=\footnotesize\sffamily,
  link/.style={-{Stealth[length=1.55mm]},line width=.73pt},
  nodebox/.style={draw=orange!39,fill=orange!17,rounded corners=3pt,minimum width=17mm,
    minimum height=9.2mm,inner sep=3pt}]
\node (input) at (0,0) {输入};
\node[nodebox] (lm) at (4.14,0) {大语言模型};
\node (output) at (8.80,0) {输出};
\node[align=center] (fixed) at (8.12,2.02) {去偏后的\\输出};
\draw[link,blue!62!black] (input.east) -- node[above] {预处理} (lm.west);
\draw[link,violet!76] (lm.east) -- node[above] {内部处理} (output.west);
\draw[link,green!45!black] (8.14,.23) -- node[right] {后处理} (fixed.south);
\draw[link,red!56] ([xshift=4.5pt]lm.south) -- ++(0,-.49) -|
  ([xshift=-4.5pt]lm.south);
\node[text=red!59!black] at (4.14,-1.22) {训练阶段};
\end{tikzpicture}

\caption{偏见缓解策略及其在大语言模型工作流程中的位置。}
\label{fig:ch06-debiasing}
\end{figure}

\subsubsection{预处理阶段的缓解}
这类方法调整模型的输入，即数据和提示，保持模型的可训练参数不变，如图~\ref{fig:ch06-debias-pre} 所示。

\sourcefigure{fig06_04.pdf}{用预处理消除偏见的方法。上半图展示数据增强、筛选、重新加权及数据生成，下半图展示改写提示文字、使用控制词元及连续提示调优。图中保留数据修改前后的英文样例。}{fig:ch06-debias-pre}

\paragraph*{数据增强}
基于数据增强的方法会将代表不足的群体的样例加入训练数据，以增加数据集的多样性。反事实数据增强（counterfactual data augmentation，CDA）是其中的重要方法。通过替换受保护属性词，例如带性别的代词，使数据集更均衡。\refcite{ch06-ref-068}{Lu 等人（2020）}以此应对职业与性别相关的偏见，在保留语法和语义完整性的同时，互换表示性别的词。

有选择地替换词是 CDA 的另一种做法，可提高数据使用效率，并针对最能发挥偏见缓解作用的样例。

\begin{practicaltip}
\refcite{ch06-ref-072}{Maudslay 等人（2019）}提出的反事实数据替换（counterfactual data substitution，CDS）等技术，将含有性别词的文本随机替换成反事实版本。同一组作者还提出人名干预（Names Intervention），专门替换名字，将男性名字与成对的女性名字关联起来再作替换。

基于 \refcite{ch06-ref-127}{Zhang 等人（2017）}的 mixup 技术，插值方法将反事实增强后的训练样例与原始版本混合，进一步丰富训练数据。
\end{practicaltip}

\paragraph*{数据筛选与重新加权}
数据筛选和重新加权可以弥补数据增强在缓解偏见方面的局限，重点是在已有数据集中精确选定要处理的样例。主要分为数据集筛选和样例重新加权两类。

数据集筛选选择数据的子集，以影响模型在微调时的学习。\refcite{ch06-ref-032}{Garimella 等人（2022）}从代表不足的群体中选取文本，增加数据多样性。\refcite{ch06-ref-010}{Borchers 等人（2022）}则选择偏见最少的样例，构造偏见较少的数据集。还有些方法针对偏见最明显的样例，通过中性化或筛除它们，减少数据集总体上的偏见。例如，\refcite{ch06-ref-107}{Thakur 等人（2023）}筛选出一组偏见明显的样例，再将与性别有关的词改为中性词，以构造更加均衡的训练数据。

样例重新加权则调整训练时特定样例的重要性。\refcite{ch06-ref-043}{Han 等人（2022）}采用这一方法，按样例标签和相关受保护属性的出现情况，给损失函数中的各个样例分配反比权重。\refcite{ch06-ref-110}{Utama 等人（2020）}和 \refcite{ch06-ref-087}{Orgad 与 Belinkov（2023）}等方法，重点降低包含社会群体信息的样例的权重，以减少模型预测时对刻板印象捷径的依赖。

\paragraph*{数据生成}
数据生成旨在解决数据增强、筛选和重新加权的局限。其中一个难点是为每个偏见维度识别具体样例，而偏见维度可能因语境或应用而异。数据生成依据预先设定的标准或特征构造全新的数据集。\refcite{ch06-ref-103}{Solaiman 和 Dennison（2021）}开发了迭代流程来构建面向特定价值观的数据集，例如消除与法律保护类别有关的偏见。撰写者在这个过程中发挥重要作用，编写能体现目标行为的提示及补全，再根据表现评估的结果改进这些样例。\refcite{ch06-ref-022}{Dinan 等人（2019）}同样请人收集多种样例，以减少聊天对话模型的性别偏见。

数据生成的一个重点是建立新词表，特别是供 CDA、CDS 等词语互换技术使用的词表。这些词表往往侧重与不同社会群体相关的词，涉及性别、种族、年龄和方言等方面。不过，依赖这类词表有时会限制所能处理的刻板印象范围。为此，有研究者提出范围更广的框架。\refcite{ch06-ref-086}{Omrani 等人（2023）}着眼于“热情程度”和“能力”等较一般的维度来理解刻板印象，扩大了缓解偏见的范围。他们为这两类特征分别生成词表，在偏见缓解任务中，可以用这些词表替代按群体构建的词表，例如性别词表。

\paragraph*{指令调优}
这里的指令调优方法会修改交给模型的输入或提示。可以为提示添加文字指令或触发内容，鼓励生成没有偏见的输出。例如，\refcite{ch06-ref-071}{Mattern 等人（2022）}用抽象程度不同的提示引导模型避开刻板印象。\refcite{ch06-ref-111}{Venkit 等人（2023）}也用对抗触发词减少国籍偏见，\refcite{ch06-ref-002}{Abid 等人（2021）}则用简短的短语减轻对穆斯林的偏见。这些方法通常把短语附加在输入上，引导模型对某个社会群体表达中性或正面的情感。

还可以用控制词元（control token）为提示分类。模型学习把每个词元与某一类输入联系起来，从而在推理时控制生成。\refcite{ch06-ref-022}{Dinan 等人（2019）}使用这一办法缓解对话生成中的性别偏见，按照训练样例中有无性别词附加对应的词元。

\begin{practicaltip}
连续提示调优（continuous prompt tuning）是在输入上加入可训练前缀的另一项技术。原预训练模型的参数保持冻结，同时允许针对具体任务调节前缀。这种办法能实现手工提示工程难以完成的、有针对性且可扩展的调整。\refcite{ch06-ref-025}{Fatemi 等人（2021）}和 \refcite{ch06-ref-120}{Yang 等人（2023）}用连续提示调优缓解性别偏见，鼓励模型使用不随受保护属性而改变的中性语言。
\end{practicaltip}

\subsubsection{训练阶段的缓解}
训练阶段的缓解策略旨在模型训练过程中减少偏见。这些技术改变优化过程，包括修改损失函数、更新下一个词的概率、在微调时有选择地冻结参数，以及移除与有害输出相关的神经元，如图~\ref{fig:ch06-debias-train} 所示。所有这些策略均通过基于梯度的训练更新来改变模型参数。

\sourcefigure{fig06_05.pdf}{训练过程中用于减少偏见的几种方法。从左到右为修改架构、修改损失函数、有选择地更新参数，以及筛除模型参数。}{fig:ch06-debias-train}

\paragraph*{修改架构}
修改架构是训练阶段缓解偏见的一个方面，涉及改变模型结构，例如层、编码器和解码器的数量、大小及类型。一个例子是加入去偏适配器模块，如 \refcite{ch06-ref-057}{Lauscher 等人（2021）}提出的 ADELE。它们以模块化适配器框架为基础，在已有层之间插入新层，以高效微调。新插入的层参与微调，预训练层保持固定，使新增部分专门学习去偏所需的信息。

Liu 等人（2022）引入一个正则项，用于尽量减小受保护属性由 $E(\cdot)$ 给出的嵌入之间的距离。
\begin{equation}
\mathcal{R}=\lambda\sum_{(a_i,a_j)\in A}\lVert E(a_i)-E(a_j)\rVert_2.
\label{eq:ch06-embedding-regularization}
\end{equation}
其中，$\mathcal{R}$ 是正则项，$\lambda$ 为缩放因子，$a_i$、$a_j$ 是集合 $A$ 的元素，分别表示受保护属性及其反事实替换。

\paragraph*{修改损失函数}
\refcite{ch06-ref-088}{Park 等人（2023）}提出一种技术，将基于投影的偏见缓解方法纳入损失函数，专门处理职业词中的性别刻板印象。他们引入一个正则项，使刻板印象相关词的嵌入 $\symbf{w}$ 和嵌入空间中的性别方向 $\symbf{g}$ 正交。这个正则项使中性职业词的嵌入远离本身具有性别含义的词，例如“sister”或“brother”。

性别方向的形式定义为
\begin{equation}
\symbf{g}=\frac{1}{|A|}\sum_{(a_i,a_j)\in A}\bigl(E(a_j)-E(a_i)\bigr),
\label{eq:ch06-gender-direction}
\end{equation}
其中，$A$ 是所有本身具有性别含义的女性相关词 $a_i$ 与男性相关词 $a_j$ 所构成的词对集合，$E(\cdot)$ 计算模型的嵌入。

正则项表示为
\begin{equation}
\mathcal{R}=\sum_{\symbf{w}\in W_{\mathrm{stereo}}}
\frac{\symbf{g}}{\lVert\symbf{g}\rVert}\symbf{w}^{\top},
\label{eq:ch06-projection-regularization}
\end{equation}
其中，$W_{\mathrm{stereo}}$ 为刻板印象相关词的嵌入集合。%
\footnote{译注：原书在公式末尾和正式编号（6.20）之间还印有“(34)”，此处只保留本书公式编号。}

\refcite{ch06-ref-004}{Attanasio 等人（2022）}考察了修改语言模型注意力层的方法。他们假设，注意力层是偏见的主要编码位置。他们提出基于熵的注意力正则化（entropy-based attention regularization，EAR），用注意力权重分布的熵衡量上下文词语的相关性。熵高，表示利用的上下文较多样。熵低，则表示模型较依赖特定词元。EAR 旨在最大化注意力权重的熵，避免对带刻板印象的词过拟合，从而让模型关注更广的输入上下文。

具体做法是在损失函数中加入以最大化熵为目标的正则项。
\begin{equation}
\mathcal{R}=-\lambda\sum_{\ell=1}^{L}\operatorname{entropy}(A)_\ell,
\label{eq:ch06-attention-entropy}
\end{equation}
其中，$\operatorname{entropy}(A)_\ell$ 是第 $\ell$ 层的注意力熵。

也有研究通过设计损失函数，使语言模型输出中与特定人口群体相关的词获得均衡的概率。例如，\refcite{ch06-ref-092}{Qian 等人（2019）}引入使概率相等的目标，鼓励模型以相同的概率预测不同人口群体的相关词。该方法有一项正则项，比较二元性别词对的 Softmax 输出概率。\refcite{ch06-ref-032}{Garimella 等人（2022）}后来将这一项用于由两类种族相关词组成的词对。对由属性 $a_i$ 和 $a_j$ 组成的 $K$ 个词对，可使用下列正则项。其中，$a_i\in A_i$ 是与社会群体 $G_i$ 有关的受保护属性词，类似地，$a_j\in A_j$ 是与社会群体 $G_j$ 有关的受保护属性词。
\begin{equation}
\mathcal{R}=\lambda\frac{1}{K}\sum_{k=1}^{K}
  \left|\log\frac{P(a_i^{(k)})}{P(a_j^{(k)})}\right|.
\label{eq:ch06-equalizing-regularization}
\end{equation}

\paragraph*{有选择地更新或筛除模型参数}
用增强的数据集微调人工智能模型可以减少偏见，但可能导致“灾难性遗忘”，使模型失去先前学到的信息。为防止这种情况，有些方法只更新极少一部分模型参数，冻结其他参数。例如，\refcite{ch06-ref-035}{Gira 等人（2022）}在 WinoBias 和 CrowS-Pairs 数据集上微调模型时，只更新层归一化等特定参数。\refcite{ch06-ref-094}{Ranaldi 等人（2023）}则只调整注意力矩阵，并冻结其他所有参数。\refcite{ch06-ref-123}{Yu 等人（2023）}有针对性地优化对某类偏见贡献最大的权重，例如与性别和职业相关的权重。

除了通过微调更新参数来减轻偏见，还有技术侧重有选择地筛除某些参数，例如在训练中或训练后把它们置零。\refcite{ch06-ref-050}{Joniak 和 Aizawa（2022）}使用移动剪枝（movement pruning），从神经网络中有选择地移除权重。他们借此从预训练模型的注意力头中选出偏见较少的一组权重。微调时冻结这些权重，并另外优化用于去偏的分数，决定移除哪些权重。

\subsubsection{推理时内部处理}
按 \refcite{ch06-ref-097}{Savani 等人（2020）}的定义，内部处理（intra-processing）在推理阶段改变预训练或微调模型的行为，无需进一步训练便可生成经过去偏的预测。这类推理阶段缓解办法，包括调整解码策略、在训练后修改模型参数，以及在推理时以模块化方式采用独立的去偏网络，如图~\ref{fig:ch06-debias-infer} 所示。

生成输出词元时，可以调整解码算法，尽量减少带有偏见的语言。这样的调整不改变可训练的模型参数。它通过选择约束、改变词元概率分布或引入辅助偏见检测模型，影响后续词语或序列出现的概率。其中一种办法是受约束的下一词元搜索（constrained next-token search），通过加入额外准则，改变输出序列中下一个词元的排序，从而引导模型给出偏见较少的输出。

\refcite{ch06-ref-033}{Gehman 等人（2020）}建议在解码时屏蔽词语或 $n$ 元组，防止选中预定义不安全词表中的词元。不过，由若干单独看并无偏见的词元或 $n$ 元组组成的文本，仍可能含有偏见，所以这一办法无法完全消除生成偏见输出的可能。

\sourcefigure{fig06_06.pdf}{推理时缓解偏见的几种方法。左侧展示受约束的下一词元搜索和改变词元概率分布，右侧展示注意力权重重分配和模块化去偏网络。}{fig:ch06-debias-infer}

\begin{practicaltip}
权重重分配（weight redistribution）是在训练后修改与已编码的偏见有关的模型权重，主要针对注意力权重。\refcite{ch06-ref-125}{Zayed 等人（2023）}通过受一个超参数控制的温度缩放，在训练后调整注意力权重。这样可以增大熵，使模型考虑更广范围的词元，可能因此避开刻板印象。也可以减小熵，使模型集中在较窄的上下文上，减少接触刻板印象词元的机会。
\end{practicaltip}

最后，模块化去偏网络解决了一些训练阶段方法的局限。训练阶段的方法通常只针对一种偏见，而且会永久改变模型。模块化去偏更灵活，可以建立独立组件，为不同任务或不同偏见将其与预训练模型组合。\refcite{ch06-ref-044}{Hauzenberger 等人（2023）}将 \refcite{ch06-ref-040}{Guo 等人（2020）}提出的差分剪枝（diff pruning）用于去偏。他们模拟同时训练多个模型，每个模型针对一个不同的维度减少偏见。再把对预训练模型参数所作的修改提取出来，高效保存在稀疏子网络中。

\subsubsection{后处理阶段的缓解}
后处理阶段的缓解对预训练模型的输出做处理，以减少偏见。如果模型主要作为黑箱提供，难以得知其训练数据或内部工作方式，这类办法尤其有用，如图~\ref{fig:ch06-debias-post} 所示。

\sourcefigure{fig06_07.pdf}{后处理阶段缓解偏见的例子。上方路线检测并替换关键词，下方路线由神经机器翻译模型改写输出。图中保留替换前后的英文句子。}{fig:ch06-debias-post}

\begin{practicaltip}
基于改写的方法检测并替换模型输出中的偏见词语或有害词语。关键词替换等技术识别带有偏见的词元，用较中性或更有代表性的用词代替，力求保留原文的内容和风格。例如，\refcite{ch06-ref-108}{Tokpo 和 Calders（2022）}利用 LIME，根据原文的潜在表示识别并替换造成偏见的词元。\refcite{ch06-ref-021}{Dhingra 等人（2023）}则利用 SHAP 识别并替换带有刻板印象的词，同时解释替换的原因。
\end{practicaltip}

另一种后处理技术使用机器翻译，把带有偏见的句子转换成中性或没有偏见的表达。多项研究采用基于规则的方法，为包含性别代词的句子构造性别中性版本（\refcite{ch06-ref-049}{Jain et al., 2021}；\refcite{ch06-ref-106}{Sun et al., 2021}）。这种办法从带有性别偏见的源句生成经去偏的对应句，构成平行数据，再训练机器翻译模型，将有性别偏见的句子转换为去偏版本。
